% The MPC slide's figure (spec, slide 3): one real episode on one time axis. Input once in the preamble, after
% vnslides.
%
% \vnmpcfigure{<width>}{<height>} draws, with the plot area <width> by <height>:
%   - the three forecast cones of talk/data/mpc_cone_{1,2,3}.csv as 68/95/99% bands, mean ± 1, 1.96 and 2.576 SD;
%   - the realized value W/W0 from talk/data/mpc_path.csv as the red line;
%   - a gold ring at each re-plan in talk/data/mpc_events.csv, with its label above the plot and a dotted leader
%     down to the ring, so no label sits on the bands.
% x is the trading day of the episode (day 0 = Oct 10, 2025); the event days are 21, 29 and 50. The x axis is titled
% with the year (review round, 2026-10-30: every axis is labeled). All files come from scripts/30_talk_data.jl.

% the forecast bands (vnband), the realized value (vnmpcred), the plot panel (vnplotpanel) and the bands' opacities
% come with the theme, from the deck palette in vnslides.sty

% \vnmpcband{<cone>}{<z>}{<opacity>}{<tag>}: the band mean ± z·sd of one cone, filled at that opacity
\newcommand{\vnmpcband}[4]{%
  \addplot[name path=vnmpc#1#4hi, draw=none, forget plot]
    table[col sep=comma, x=day, y expr={\thisrow{mean} + #2*\thisrow{sd}}] {data/mpc_cone_#1.csv};
  \addplot[name path=vnmpc#1#4lo, draw=none, forget plot]
    table[col sep=comma, x=day, y expr={\thisrow{mean} - #2*\thisrow{sd}}] {data/mpc_cone_#1.csv};
  \addplot[fill=vnband, fill opacity=#3, draw=none, forget plot] fill between[of=vnmpc#1#4hi and vnmpc#1#4lo];
}

% \vnmpccone{<cone>}: the 99, 95 and 68% bands of one cone, widest first so the inner bands layer on top, each at the
% theme's opacity for it
\newcommand{\vnmpccone}[1]{%
  \vnmpcband{#1}{2.576}{\vnbandouter}{a}%
  \vnmpcband{#1}{1.96}{\vnbandmiddle}{b}%
  \vnmpcband{#1}{1.0}{\vnbandinner}{c}%
}

% \vnmpcevent{<row>}{<node options>}{<label>}: the ring at the re-plan in row <row> (from 0) of mpc_events.csv, on the
% red line, a dotted leader from the top of the plot down to it, and its label above the plot, placed by <node
% options> (anchor=south east, align=right ends the label at the leader; anchor=south west, align=left starts it
% there). The path has one row per day from day 0, so the value on the event day is the path's row of that number.
% \\ breaks the label. pgfplots runs the drawing commands at \end{axis}, so \edef fixes the day and value now, and
% \unexpanded carries the label's own tokens.
\newcommand{\vnmpcevent}[3]{%
  \pgfplotstablegetelem{#1}{day}\of\vnmpcEvents
  \let\vnmpcDay\pgfplotsretval
  \pgfplotstablegetelem{\vnmpcDay}{w}\of\vnmpcPath
  \let\vnmpcW\pgfplotsretval
  \edef\vnmpcCmd{%
    \noexpand\draw[vnfigaccent, line width=0.6pt, dash pattern=on 1pt off 1.5pt]
      (axis cs:\vnmpcDay,\vnmpcYtop) -- (axis cs:\vnmpcDay,\vnmpcW);
    \noexpand\node[circle, draw=vnfigaccent, line width=1.2pt, fill=vnplotpanel, inner sep=0pt, minimum size=5pt]
      at (axis cs:\vnmpcDay,\vnmpcW) {};
    \noexpand\node[#2, font=\noexpand\fontsize{7.5}{9}\noexpand\selectfont, text=vnslideink, inner sep=0pt,
      yshift=3pt]
      at (axis cs:\vnmpcDay,\vnmpcYtop) {\unexpanded{#3}};}%
  \vnmpcCmd
}

\newcommand{\vnmpcfigure}[2]{%
  \pgfplotstableread[col sep=comma]{data/mpc_path.csv}\vnmpcPath
  \pgfplotstableread[col sep=comma]{data/mpc_events.csv}\vnmpcEvents
  % the band labels sit at the upper edge of each band where the last cone ends
  \pgfplotstableread[col sep=comma]{data/mpc_cone_3.csv}\vnmpcLast
  \pgfplotstablegetrowsof{\vnmpcLast}%
  \pgfmathtruncatemacro\vnmpcLastRow{\pgfplotsretval - 1}%
  \pgfplotstablegetelem{\vnmpcLastRow}{day}\of\vnmpcLast
  \let\vnmpcEndDay\pgfplotsretval
  \pgfplotstablegetelem{\vnmpcLastRow}{mean}\of\vnmpcLast
  \let\vnmpcEndMean\pgfplotsretval
  \pgfplotstablegetelem{\vnmpcLastRow}{sd}\of\vnmpcLast
  \let\vnmpcEndSd\pgfplotsretval
  % how far below its forecast the band exit was, in SDs, from the band_exit row (row 1; the data script checks the
  % order of the re-plans): the exit is too close to the band edge to see, so the label states it
  \pgfplotstablegetelem{1}{z}\of\vnmpcEvents
  \pgfmathparse{abs(\pgfplotsretval)}%
  \pgfmathprintnumberto[fixed, precision=1, assume math mode=true]{\pgfmathresult}{\vnmpcExitZ}%
  % the top of the plot, where the re-plan labels hang
  \def\vnmpcYtop{1.25}%
  \begin{tikzpicture}
    \begin{axis}[
      width=#1, height=#2, scale only axis,
      clip=false,
      xmin=0, xmax=50, ymin=0.90, ymax=\vnmpcYtop,
      axis background/.style={fill=vnplotpanel},
      axis line style={draw=none},
      grid=major, major grid style={draw=vnplotgrid, line width=0.8pt},
      tick style={draw=none},
      xtick={0,21,29,50}, xticklabels={Oct 10, Nov 10, Nov 20, Dec 22},
      % no label at the bottom edge, 0.90, so the first date has the corner to itself
      ytick={0.95,1.00,1.05,1.10,1.15,1.20,1.25},
      tick label style={font=\fontsize{9}{10.8}\selectfont, color=vndottick},
      yticklabel style={/pgf/number format/assume math mode=true, /pgf/number format/fixed,
        /pgf/number format/fixed zerofill, /pgf/number format/precision=2},
      ylabel={$W/W_0$},
      xlabel={date, 2025}, xlabel style={yshift=3pt},
      label style={font=\fontsize{9.5}{11.4}\selectfont, color=vnslidebody},
    ]
      % forecast cones, each from the value at its re-plan
      \vnmpccone{1}%
      \vnmpccone{2}%
      \vnmpccone{3}%

      % band labels right of the last cone's end
      \pgfmathsetmacro\vnmpcYa{\vnmpcEndMean + 2.576*\vnmpcEndSd}%
      \pgfmathsetmacro\vnmpcYb{\vnmpcEndMean + 1.96*\vnmpcEndSd}%
      \pgfmathsetmacro\vnmpcYc{\vnmpcEndMean + 1.0*\vnmpcEndSd}%
      \node[anchor=west, xshift=3pt, font=\fontsize{8}{9.6}\selectfont, text=vnslidebody]
        at (axis cs:\vnmpcEndDay,\vnmpcYa) {99\%};
      \node[anchor=west, xshift=3pt, font=\fontsize{8}{9.6}\selectfont, text=vnslidebody]
        at (axis cs:\vnmpcEndDay,\vnmpcYb) {95\%};
      \node[anchor=west, xshift=3pt, font=\fontsize{8}{9.6}\selectfont, text=vnslidebody]
        at (axis cs:\vnmpcEndDay,\vnmpcYc) {68\%};

      % the realized value
      \addplot[color=vnmpcred, line width=1.8pt, line join=round, mark=none]
        table[col sep=comma, x=day, y=w] {data/mpc_path.csv};

      % the re-plans, labelled above the plot: the first and last labels end at their leaders, the band exit's starts
      % at its leader, so the two close together (8 days apart) do not collide
      \vnmpcevent{0}{anchor=south east, align=right, xshift=-2pt}{{\vnmedium Re-plan}\\21 days passed}%
      \vnmpcevent{1}{anchor=south west, align=left, xshift=2pt}{{\vnmedium Re-plan}\\left the 95\% band\\
        \vnmpcExitZ{} SD below the mean}%
      \vnmpcevent{2}{anchor=south east, align=right, xshift=-2pt}{{\vnmedium Re-plan}\\21 days passed}%
    \end{axis}
  \end{tikzpicture}%
}
